% Independent mathematical audit, 2026-11-14.
% This is review material, not a verified proof or a substitute for elaborated Lean interfaces.
\section{Rectifiability branch: proof sketch and interface audit}

This audit uses the supplied De Lellis PDF: Chapter 4 is printed pp.~27--39
(PDF pp.~31--43), Chapter 5 is printed pp.~40--55 (PDF pp.~44--59).
References below refer to this version, not the older 2006 lecture-note pagination.
The input document supplies mathematics, not workflow instructions.
The published author errata must also be consulted:
\url{https://www.math.ias.edu/delellis/sites/math.ias.edu.delellis/files/Errata_libro_0.pdf}.

\subsection{What the main Preiss proof actually needs}

Let $H=\mathcal H^m$ with the project's fixed normalization, and let
$D_{x,r}\mu=r^{-m}(T_{x,r})_\#\mu$.
For $0<m\leq n$, the rectifiability input to Preiss is the implication
\[
\left[0<\Theta_*^m(\mu,x)\leq\Theta^{*m}(\mu,x)<\infty
\ \hbox{and every tangent is flat, for $\mu$-a.e. $x$}\right]
\Longrightarrow \mu\text{ is }m\text{-rectifiable}.
\]
This is Theorem 5.1. It does not assume existence of the density limit,
and does not assume that all tangent measures at a point use the same plane.
The latter stronger property follows almost everywhere after rectifiability,
using Chapter 4. It must not be substituted for the hypothesis in the main proof.

Here is the complete architecture of the Chapter 5 argument.
First stratify the positive finite density bounds on bounded Borel pieces.
Covering comparisons show that these pieces have finite $H$ measure and that
$\mu$ and $H$ restricted to each piece are mutually absolutely continuous,
with uniformly positive finite comparison constants after further stratification.
Use Radon--Nikodym on each finite restriction, never on ambient $H$ when $m<n$.
Differentiate the density and restrictions to transfer the hypotheses to
$\nu=H|F$ on each finite piece. Consequently, $\nu$ has positive lower and
finite upper $m$-density and every nonzero tangent is flat at almost every
point of $F$.

These two properties imply weak linear approximation at each such point.
If either clause of Definition 5.4 fails, choose bad radii decreasing to zero
with error threshold $\eta>0$ and lower-mass threshold $1/j$.
The upper density bound gives uniform mass bounds on compact sets for
$D_{x,r_j}\nu$. Compactness supplies a vaguely convergent subsequence, and
the lower density bound makes its limit nonzero. It is a tangent, hence
$cH|V$ with $c>0$. Failure of the tube clause leaves positive mass in a
compact set separated from $V$, contradicting compact-set portmanteau.
Failure of the plane-centred-ball clause gives centres in the unit disk of
$V$ with masses tending to zero. Extract a convergent subsequence of centres
and put a fixed smaller ball around the limiting centre inside the moving
balls. Open-set portmanteau contradicts the positive flat mass. This last
fixed-smaller-ball argument is the rigorous version of the informal moving-ball
limit in equation (5.36). No continuity-set theorem for moving balls is needed.

It remains to prove Proposition 5.5: a finite-Hausdorff-measure Borel set
which is weakly linearly approximable almost everywhere is rectifiable.
The proof has substantially more content than the current single graph node.
Decompose a counterexample into rectifiable and purely unrectifiable parts.
Weak linear approximation passes to the latter at density-one points,
using finite upper density to convert relative discarded mass into $o(r^m)$.
Thus a positive-measure purely unrectifiable counterexample exists.

Lemma 5.9 then proves that every orthogonal projection of this counterexample
onto an $m$-plane has zero $H$ measure. This is stronger than the usual
almost-every-plane projection theorem and uses weak linear approximation.
Uniformize on large compact subsets: lower mass bounds hold uniformly, and
weak tube concentration upgrades to actual containment of the compact subset
in a thin tube. A point of the compact subset outside the tube would carry a
small ball with too much mass. Pure unrectifiability and the cone-to-graph
lemma imply that almost every point has arbitrarily close neighbours almost
orthogonal to any fixed projection plane. At the scale of such a pair, the
approximating plane therefore has a unit direction with small projection.
Elementary linear algebra puts the projected tube into a thin slab. Its
$m$-measure is small. A disjoint almost-everywhere ball covering and the lower
mass bound sum these small projection estimates. The compact complement has
small measure and its projection has no larger measure. Letting that error
tend to zero proves the projection is null.

Finally, null projections are incompatible with weak linear approximation
and positive finite $H$ measure. Choose nested positive-measure compact sets
$G\subset F\subset E$, uniform lower mass bounds on $F$, and at each point of
$G$ scale-dependent approximating planes. Both directions of approximation
are needed: points of $F$ lie near the plane, and points of the plane lie near
$F$. At a good point of $G$, $H(E\cap B_r)$ has a uniform upper bound and
$H((E\setminus G)\cap B_{2r})$ is as small as desired compared with $r^m$.
Fix its approximating plane $W$ at scale $\sigma$.
Since $P_W(G)$ is null, almost all of the $W$ disk consists of holes in the
projection. Select disjoint disks at their distance-to-projection radii.
The sum of their $m$th powers is bounded below by a constant times $\sigma^m$.
Discard disks whose inner cylinders meet $F$: the discarded radius sum is
controlled by the small mass of $E\setminus G$.
For each remaining disk, a nearest projected point of compact $G$ gives a
boundary point $z_i$. The approximating plane there must contain a direction
almost perpendicular to $W$; otherwise it enters the empty inner cylinder.
Along this direction place at least a constant times $1/\eta$ separated balls.
Plane-to-$F$ approximation moves their centres onto $F$, so every ball has a
uniform lower mass bound. All these balls, across different disks as well,
are disjoint. Their total mass is at least
$c\delta\sigma^m/\eta$, contradicting the upper mass bound when $\eta$ is
chosen sufficiently small. This proves Proposition 5.5 and then Theorem 5.1.

\subsection{Definitions and theorem statements which must be explicit}

The following inventory supplies boundaries for reusable Lean declarations.
Each mathematical bullet should have its own declaration or a small API,
not merely a sentence hidden inside the proof of the final criterion.
Definitions should not build in local finiteness or measurability unless
those properties are intrinsic; theorems should state such hypotheses.

\begin{enumerate}
\item \textbf{Cones and graph parametrizations.}
For a finite-dimensional plane $V$, use the existing orthogonal projection
$P_V$, define $Q_V=\mathrm{id}-P_V$, and define
$C(x,V,\alpha)=\{y:\|Q_V(y-x)\|\leq\alpha\|P_V(y-x)\|\}$.
Record translations, positive homotheties, closedness, and aperture monotonicity.
For a Lipschitz $f:V\to V^\perp$, the graph is the range of $v\mapsto v+f(v)$.
Record that this map is Lipschitz and antilipschitz, its inverse on the graph
is $P_V$, the full graph is closed, and bounded graph pieces have finite $H$ measure.

\item \textbf{Cone-to-graph, De Lellis Lemma 4.7.}
For an arbitrary set $F$, if $y\in C(x,V,\alpha)$ for all $x,y\in F$,
with $\alpha\geq0$, then $F$ is contained in a full Lipschitz graph over $V$.
No measure hypotheses are needed. The graphing function first exists on
$P_V(F)$; extend into $V^\perp$ using finite-dimensional coordinates.
The empty-set case should avoid choosing an element of $F$.

\item \textbf{Countable rectifiability calculus.}
Empty sets, subsets, countable unions, Lipschitz images, isometric images,
Hausdorff-null modifications, Borel carriers, and sigma-finite Hausdorff
restriction to a countably parametrized Euclidean carrier need separate lemmas.
The Borel carrier is the union of compact images of closed bounded domain
pieces after extending parametrizations into a Euclidean target.

\item \textbf{Pure unrectifiability, Definition 5.6.}
For the canonical reusable definition require $H(E\cap R)=0$ for every
countably rectifiable $R$. This directly gives heredity and the null-intersection
API, and it implies the graph-test condition of De Lellis Definition 5.6
because every graph is rectifiable. Equivalence with graph tests can be proved
on the companion image-to-graph branch. Finite measure belongs in later
theorem hypotheses, not in the definition. This choice keeps the core
decomposition and Chapter 5 argument independent of the area-formula branch.

\item \textbf{Finite-measure decomposition, Theorem 5.7.}
If $E$ is Borel and $H(E)<\infty$, obtain Borel disjoint subsets $R,U$ with
$E=R\cup U$, $R$ rectifiable, and $U$ purely unrectifiable. Prove uniqueness
modulo $H$-null sets. Use a supremum and a countable maximizing family, then
rectifiability of a countable union. This does not require the area formula.

\item \textbf{Density-to-Hausdorff comparisons.}
State local comparison inequalities on sets where all sufficiently small
ball masses obey common power bounds. Derive absolute continuity from finite
upper density; derive finite Hausdorff measure of bounded lower-bound strata;
then obtain countably many Borel finite-measure pieces and two-sided measure
comparisons. This is the reusable content behind Proposition 2.16 and the
current bundled finite-pieces node. Lower density may be weakened in other
versions, but the main Preiss route already supplies the stronger hypothesis.

\item \textbf{Density and tangent transfer to pieces.}
On each finite piece use its Hausdorff restriction as the reference measure,
not ambient $H$. Establish positivity/finiteness of the Radon--Nikodym density
at the relevant almost-everywhere points, its local integrability, and
restriction regularity. Transfer arbitrary-scalar tangent sets using GMT.
For fixed normalization, additionally prove a compact-test difference tends
to zero with the explicit density multiplier. Mere equality of GMT tangent
sets loses this multiplier.

\item \textbf{Fixed-normalization compactness at a density-bound point.}
Given $0<\Theta_*^m(\nu,x)\leq\Theta^{*m}(\nu,x)<\infty$ and any positive
$r_j\to0$, obtain a subsequence of $D_{x,r_j}\nu$ converging to a nonzero
regular measure. Obtain local upper bounds from density and bound the finite
initial segment separately. A closed-ball lower mass bound and portmanteau
prove nonzeroness. Show the resulting limit satisfies GMT's tangent predicate.

\item \textbf{Flat tangents imply weak linear approximation.}
For $H(F)<\infty$, at a point with the preceding density bounds where every
GMT tangent of $H|F$ is $m$-flat, prove the two clauses of Definition 5.4.
The theorem is pointwise and should not repeat the almost-everywhere
localization to $F$. Separate compact-tube and moving-ball contradiction lemmas.

\item \textbf{Weak approximation gives positive lower density.}
Take the plane-centred-ball clause at $z=x$ and one fixed aperture. This
supplies the lower density needed for subsequent uniform strata. Do not
mistake the general Hausdorff lower bound on \emph{upper} density for this result.

\item \textbf{Weak approximation survives negligible deletion.}
If $F\subset E$ and $H((E\setminus F)\cap B_{2r}(x))=o(r^m)$, weak linear
approximation of $E$ at $x$ implies that of $F$. Shrink the lower-ball
constant and the error aperture. Combine differentiation of a restriction
with finite upper density to produce this hypothesis almost everywhere in $F$.
This is the substantive content of Lemma 5.8, beyond decomposition.

\item \textbf{Closed uniform good-set exhaustion.}
Fix a finite regular measure $\nu$, an aperture $\eta>0$, mass constant
$\lambda>0$, and radius cutoff $R>0$. Define the relaxed good set by requiring,
for every $0<r<R$, some $m$-plane $V$ with
\[
\nu\{y:|y-x|<r,\ \operatorname{dist}(y-x,V)>\eta r\}
\leq\eta r^m,
\]
and
\[
\nu(\overline B_{\eta r}(z))\geq\lambda r^m
\quad\text{for all }z\in(x+V)\cap B_r(x).
\]
This good set is closed in $x$, after giving the planes their compact
orthogonal-projection topology. For fixed $r$, the first condition is closed
in $(x,V)$ by lower semicontinuity of the mass of a moving open set. The
second is closed by upper semicontinuity of closed-ball mass and approximation
of each test centre $z$ by $x_i+P_{V_i}(z-x)$; the strict condition $|z-x|<r$
persists for large $i$. Projection of this closed relation along compact
plane space is closed, and intersection over all $0<r<R$ remains closed.
Thus no reduction to rational radii or general measurable-projection theorem
is needed for this construction. Existing weak approximation at aperture
$\eta/2$ implies a relaxed $\eta$ condition. Conversely, a relaxed
$\eta/2$ condition implies the original aperture-$\eta$ condition: enlarge
the tube and use $\overline B_{\eta r/2}\subset B_{\eta r}$, retaining the
positive lower constant and obtaining strict upper error from $\eta/2<\eta$.
A countable choice of rational lower constants and radius cutoffs now gives
Borel exhaustions after intersection with Borel $E$. Apply inner regularity
there to choose compact pieces. A measurable hull of an arbitrary good set
alone would not justify choosing compact subsets of that original set.
This compact-plane and semicontinuity API is therefore foundational for
Chapter 5 as well as for the later flatness functional.

\item \textbf{Tube concentration to support containment.}
Suppose every $z\in C$ has $H(E\cap B_s(z))\geq\delta s^m$ for small $s$.
If the mass of $E\cap B_{2r}(a)$ outside the $(\eta r/2)$-tube about $a+W$
is less than $\delta(\eta r/2)^m$, then
$C\cap B_r(a)$ lies in the $\eta r$-tube. This is a direct geometric
ball-inclusion argument and works for measures more generally than $H$.

\item \textbf{Purely unrectifiable sets escape every fixed cone.}
For every plane $V$ and fixed finite aperture, at $H$-almost every point of a
purely unrectifiable set $C$, there are points of $C$ outside that cone at
arbitrarily small distances. The exceptional set is a countable union of
local cone-condition pieces; cover each by countably many small balls and
apply the cone-to-graph lemma. For an arbitrary (possibly nonmeasurable) $C$,
use the fixed-radius bad-anchor sets
$B_r=\{a:\forall b\in C,\ d(a,b)<r\Rightarrow b\in X(a,V,A)\}$.
These sets are closed, since each implication is a union of two closed
conditions. The graph argument proves $H(C\cap B_r)=0$; measurability of
$B_r$ then gives $(H|_C)(B_r)=0$ without assuming measurability of $C$.
A countable sequence of radii completes the almost-everywhere assertion.
This is the Step 2 input to Lemma 5.9.

\item \textbf{Small projection direction gives a thin projected slab.}
For $m$-planes $V,W$, a unit $w\in W$ with $\|P_Vw\|\leq\varepsilon$
implies that the projection onto $V$ of the $\eta R$-tube about $a+W$ inside
$B_R(a)$ has $H$ measure at most $C_m(\eta+\varepsilon)R^m$.
Choose a unit vector in $V$ orthogonal to $P_V(W\cap w^\perp)$, contain the
projection in a slab, and estimate its volume by orthogonal coordinates.
Closed bounded tube variants avoid measurability problems for image sets.

\item \textbf{Every-plane null projections, Lemma 5.9.}
A Borel purely unrectifiable finite-$H$ set weakly linearly approximable
almost everywhere has null projection onto \emph{each} $m$-plane.
This needs the preceding compact uniformization, tube containment,
cone escape, slab bound, Lipschitz Hausdorff distortion, and an actual
almost-everywhere disjoint Besicovitch covering.
It is not a consequence of GMT's fivefold-covering statement alone without
additional work accounting for uncovered mass.

\item \textbf{Nested compact two-sided approximation.}
For a positive finite-$H$ weakly approximable set, extract compact
$G\subset F\subset E$ of positive measure, uniform lower mass at $F$, and
planes at points of $G$ giving both tube containment of $F$ and proximity of
plane points to $F$. Retain the lower mass condition on plane-centred balls
of $E$. Density differentiation of $F$ in $E$ upgrades that mass condition
to proximity. This is Section 5.2 Step 1, equations (5.15)--(5.21).

\item \textbf{Good base point with small discarded mass.}
Select $a\in G$ with a uniform upper bound on $H(E\cap B_r(a))$ and
$H((E\setminus G)\cap B_{2r}(a))\leq t r^m$ for all sufficiently small $r$.
Use the upper Hausdorff-density bound for $G$, and the outside-density-zero
bound for $E\setminus G$, or differentiation plus finite upper density.
Constants depend on the chosen Hausdorff normalization.

\item \textbf{Distance-to-projection disk selection and radius sum.}
For compact $G$ and fixed $W$, let the hole radius be distance to the compact
projection of $G\cap\overline B_{2\sigma}(a)$.
Its positivity on the complement, its upper bound $2\eta\sigma$, and a
fivefold covering give disjoint disks $D_{4\rho_i}$ whose $D_{20\rho_i}$
cover almost all a fixed $W$ disk. Therefore $\sum_i\rho_i^m\geq c_m\sigma^m$.
Use the closed ball in the compact projection; an intersection with an open
ball need not be compact and nearest projected points need not exist.
The compact replacement also requires a scale buffer: obtain the base tube
bound at radius $2\sigma$, so it controls $F\cap B_{4\sigma}(a)$ and hence
all of $G\cap\overline B_{2\sigma}(a)$. Write $\beta=2\eta$ for the
working aperture. Then $\rho_i\leq\beta\sigma$, and a nearest projected
witness $z$ satisfies
$\|z-a\|\leq\sigma/4+2\beta\sigma<\sigma/2$ when $\beta<1/8$.
Thus no boundary point is being subjected to an open-ball hypothesis.
Use the local radius $r_i=\rho_i/(8\beta)=\rho_i/(16\eta)$;
the original approximation error $\eta r_i=\rho_i/16$ is smaller than
the $\rho_i/8$ tolerance needed for the column construction.
All subsequent column estimates use the working aperture $\beta$;
choose $\sigma$ sufficiently small that $2\sigma<R$.

\item \textbf{Discarded cylinders have small radius sum.}
Cylinders whose inner halves meet $F$ contain small balls based at $F$ and
avoiding $G$. Lower mass and disjointness give
$\sum_{i\in J}\rho_i^m\leq C_m t\sigma^m/\delta$.
Choose $t$ after $\delta,\eta,\gamma$ so that the remaining sum retains a
positive dimensional multiple of $\sigma^m$.

\item \textbf{Empty-cylinder linear geometry.}
Let $V,W$ be $m$-planes, $p\in V$, and $\|P_Vz-p\|=\rho>0$.
If $(z+W)\cap B_r(z)$ avoids the cylinder
$\{q:\|P_Vq-p\|<\rho/4\}$ and $r$ is sufficiently larger than $\rho$,
then $W$ has a unit vector $w$ with $\|P_Vw\|\leq C\rho/r$.
If the restricted projection has kernel, this is immediate; otherwise use
its least singular value and equality of the two plane dimensions.
A segment of length $c r$ based at $z$ therefore remains in the larger
$2\rho$ cylinder. This standalone lemma supplies the geometric step omitted
as an ``easy computation'' in Section 5.2.

\item \textbf{Disjoint column packing and mass contradiction.}
At the local scale $r_i=\rho_i/(8\beta)$ with $\beta=2\eta$, the preceding segment contains at
least $c/\eta$ suitably separated centres. Approximate them by points of $F$
to distance $\rho_i/8$ and use balls of radius $\rho_i/8$.
Their masses are at least $c_m\delta\rho_i^m$, and their locations are in
pairwise disjoint $4\rho_i$ cylinders and inside $B_\sigma(a)$.
Conclude $H(E\cap B_\sigma(a))\geq c_m\delta\sigma^m/\eta$, contradicting
the upper bound. Constants and the integer rounding in the number of balls
must be stated explicitly or hidden behind proved positive-dimensional constants.

\item \textbf{Marstrand--Mattila assembly.}
The preceding contradiction proves Proposition 5.5. Apply it to each finite
Hausdorff piece, then reassemble the countable rectifiable carrier to prove
Theorem 5.1. Separately combine it with the necessity direction below for
the user's requested if-and-only-if characterization.
\end{enumerate}

\subsection{The reusable reverse direction: not a prerequisite of Preiss}

The current new statements correctly distinguish a common tangent plane
from arbitrary planes depending on individual tangents. For locally finite
rectifiable $\mu$, one has, at almost every $x$, a single $V_x$ and $c_x>0$
with convergence along every positive sequence $r_j\to0$:
\[
D_{x,r_j}\mu\rightharpoonup c_xH|V_x.
\]
Consequently, GMT's tangent set is the whole positive ray
$\{cH|V_x:0<c<\infty\}$, not a singleton. The diameter-normalized density
is $c_x H(V_x\cap\overline B_1)/(2^m)$; identify it with $c_x$ only
after proving the plane ball normalization for the chosen Hausdorff convention.

To prove this branch without circular use of Preiss, expose these obligations:
\begin{enumerate}
\item Lipschitz maps send Lebesgue-null domain sets to $H$-null images.
This follows from the Lipschitz Hausdorff distortion bound and the domain
Hausdorff/Lebesgue normalization.
\item The image of the derivative-rank-deficient part is $H$-null
(De Lellis Lemma 4.11); the full-rank differentiability locus can be partitioned
into countably many graph pieces. These prove image/graph equivalence without
requiring the entire multiplicity area formula.
\item For injective graph parametrizations, prove the affine Jacobian
change-of-variables case and then the graph area formula by differentiability
and countable derivative-control pieces (Sections 4.1--4.2).
\item At points of differentiability and Lebesgue differentiation of the
Jacobian, compactly supported integrals of the blown-up graph measure tend
to those on the derivative plane. This is a graph-level theorem.
\item Restrict to Borel graph pieces at density-one points, transfer by the
Radon--Nikodym density with its exact coefficient, and patch the countable
family. This proves Corollary 4.4 and the necessity direction of Theorem 4.8.
\item Continuity balls of the flat limit give positive finite density.
The tangent-normalization adapter identifies the complete positive ray.
Combine with Theorem 5.1 for the full characterization using two-sided bounds.
\end{enumerate}
Whitney extension and $C^1$ graph equivalence are not needed for either the
main Preiss route or this Lipschitz graph necessity proof.

\subsection{Specific source errors and cautions}

The supplied PDF was visually checked at PDF pp.~56--57. The proof really
prints $c_1\sigma=c_1\rho_i/(8\eta)$ and
$N\geq c_2\sigma/\rho_i=c_2/(8\eta)$, although only
$\rho_i\leq\eta\sigma$ was established. These equalities cannot be used.
Work at the local scale $r_i=\rho_i/(8\eta)$, obtain a segment of length
comparable to $r_i$, and derive $N\geq c/\eta$ directly. This is sufficient
for the final contradiction. The local approximation statement applies
inside $B_{r_i}(z_i)$, not automatically throughout the much larger ambient
ball displayed in the prose. The standalone empty-cylinder lemma above is
intended to repair this logical boundary; its proof must be checked.

Other typographical errors should not be copied into Lean: (5.5) lists
self-rectifiable intersections where the cross intersections are intended;
Theorem 5.7 is called Lemma 5.8 in the proof of Lemma 5.8; the lower-bound
ball in the proof following (5.10) is centred at the outlying point $z$, not
at $a$; Section 5.3 writes $c\in]1,0[$ instead of $0<c<1$; and some factors
$\omega_m$ are suppressed in the last display. The normalization-free
statements above absorb fixed positive constants correctly.

The author errata adds substantial corrections outside Chapters 4--5 which
also affect this project's whole DAG: the hyperplane in Lemma 3.8 need not
be the boundary of the original half-space; the normalized tangent-transfer
Proposition 3.12 requires finite upper density; the $m=2$ argument for (8.21)
must be replaced using the trace identities and (8.22); and measurability of
the exceptional set in the recentering argument is a nontrivial measurable
projection problem. These corrections must be separate review obligations.

\subsection{Pinned-source reuse audit}

The following names were checked in the existing checkout, not guessed.
They are starting points for imports and adapters, not a declaration that
all desired composite statements are already formalized.
\begin{itemize}
\item \texttt{Mathlib/Topology/MetricSpace/Lipschitz.lean}:
\texttt{LipschitzOnWith.extend\_real} and
\texttt{LipschitzOnWith.extend\_pi}; transport finite coordinate extensions
back to Euclidean subspaces. The optimal Kirszbraun constant is unnecessary.
\item \texttt{Mathlib/Analysis/Calculus/Rademacher.lean}:
\texttt{LipschitzWith.ae\_differentiableAt}.
\item \texttt{Mathlib/MeasureTheory/Measure/Hausdorff.lean}:
\texttt{LipschitzOnWith.hausdorffMeasure\_image\_le},
\texttt{LipschitzWith.hausdorffMeasure\_image\_le},
\texttt{Isometry.hausdorffMeasure\_image},
\texttt{Isometry.map\_hausdorffMeasure}, homothety scaling, and
\texttt{MeasureTheory.hausdorffMeasure\_orthogonalProjectionOnto\_le}.
Beware: \texttt{hausdorffMeasure\_pi\_real} uses the function-space norm,
not automatically the Euclidean $\ell^2$ norm.
\item \texttt{Mathlib/MeasureTheory/Covering/Besicovitch.lean}:
\texttt{Besicovitch.exists\_disjoint\_closedBall\_covering\_ae},
\texttt{Besicovitch.ae\_tendsto\_rnDeriv}, and
\texttt{Besicovitch.ae\_tendsto\_measure\_inter\_div}.
The density-ratio theorem needs both measures locally finite, satisfied on
finite Hausdorff pieces. Its balls are closed; supply open/closed adapters.
\item \texttt{Mathlib/MeasureTheory/Covering/Vitali.lean}:
\texttt{Vitali.exists\_disjoint\_subfamily\_covering\_enlargement\_ball}
for the fivefold disk selection. It is distinct from an almost-everywhere
covering by the disjoint balls themselves.
\item \texttt{Mathlib/MeasureTheory/Measure/Decomposition/}:
Radon--Nikodym, \texttt{Measure.rnDeriv\_ne\_top}, measurable density
representatives, and with-density absolute continuity. Instantiate these on
$H|F$ or a sigma-finite Borel carrier, not on ambient $H$.
\item FormalizingGMT's existing graph nodes supply finite-restriction
regularity, upper-density estimates inside and outside finite Hausdorff sets,
vague compactness, portmanteau, blow-up ball identities, and tangent matching.
\item Further reusable declarations in
\texttt{FormalizingGMT/TangentMeasures/MatchingTangents.lean} are
\texttt{tendsto\_nhdsGT\_zero\_of\_pos},
\texttt{norm\_integral\_cc\_le},
\texttt{weaklyConverges\_of\_tendsto\_integral\_sub},
\texttt{eventually\_smul\_measure\_closedBall\_le}, and
\texttt{tendsto\_mul\_measure\_diff\_zero}. These are useful for the
fixed-dimensional-normalization restriction and density adapters too.
\end{itemize}

\subsection{Recommended file boundaries and genericity}

Use dependency-driven files rather than one monolithic statement laboratory.
A project namespace is suitable while developing, but reusable concepts
should have mathlib-style logical ownership and minimal imports.
\begin{itemize}
\item \texttt{Geometry/Planes}: finite-dimensional linear planes and
orthogonal projection geometry, independent of measures and tangent predicates.
A Euclidean space abbreviation should live in a minimal geometry/basic file,
not require importing rectifiability definitions.
\item \texttt{Rectifiability/Defs}, \texttt{Rectifiability/Basic}:
canonical countably rectifiable sets, rectifiable measures, countable-carrier
calculus, and almost-everywhere invariance.
\item \texttt{Rectifiability/Graphs}, \texttt{Rectifiability/Cones}:
graph parametrization API, coordinate extensions, cone-to-graph geometry.
\item \texttt{Rectifiability/PurelyUnrectifiable} and
\texttt{Rectifiability/Decomposition}: the pure predicate and finite-measure
decomposition, using graph coverage but not tangent measures.
\item \texttt{Density/HausdorffComparison}, \texttt{Density/FinitePieces}:
density-to-Hausdorff comparisons and sigma-finite carrier stratification.
\item \texttt{Rectifiability/DensityRepresentation}:
Radon--Nikodym characterization and bounded-density pieces.
\item \texttt{Tangent/Restriction}, \texttt{Tangent/DensityTransfer}:
adapters preserving dimensional normalization and its exact coefficient.
\item \texttt{Rectifiability/WeakApproximation/Defs},
\texttt{.../Restriction}, \texttt{.../Uniformization}:
the two-clause approximation predicate, deletion stability and closed good-set
uniform compact selection. These definitions only need geometry and
Hausdorff measure; do not import tangent or uniform-measure theory here.
\item \texttt{Rectifiability/ProjectionEstimate},
\texttt{Rectifiability/MarstrandMattila/ProjectionNull},
\texttt{.../Packing}, \texttt{.../SetCriterion}, \texttt{.../MeasureCriterion}:
isolate the long Chapter 5 chain and reusable linear geometry from its
measure-theoretic contradiction.
\item \texttt{Rectifiability/GraphCover}, \texttt{Rectifiability/GraphArea},
\texttt{Rectifiability/GraphBlowup}, \texttt{Rectifiability/TangentCharacterization}:
the reverse-direction branch and full characterization. It should not enter
the import closure of the Preiss assembly unless actually used there.
\end{itemize}

For arbitrary metric targets, the canonical definition should allow Lipschitz
maps on subsets of Euclidean space. A globally parametrized coverage predicate
is useful separately, and agrees with the standard definition for Euclidean
targets after an extension theorem. It is too restrictive in general: a
positive-measure fat Cantor set, considered as a metric target, is parametrized
by its identity on that subset of the line, while every continuous map from
the whole line into it is constant. Subtype-domain maps avoid empty-target
extension issues. One can define rectifiable measures by absolute continuity
with respect to $H$ and concentration on a countably rectifiable carrier;
absolute continuity absorbs the carrier's Hausdorff-null remainder.

The current density-representation equivalence is mathematically sound with
its explicit Borel and sigma-finite carrier. Its converse remains valid for
nonnegative extended-valued densities: a Hausdorff-null remainder has zero
weighted measure even when the weight takes the value infinity there.

This audit does not certify the current Lean source as elaborated. In particular,
all technical nodes listed above still require exact checked Lean interfaces
and proofs; source availability and a rendered DAG are not proof verification.

The machine-readable companion to this review is \texttt{rectifiability-nodes.json}:
53 proposed interface nodes, comprising 42 core and 11 companion nodes.
Existing labels are refined rather than duplicated. In particular, move
compact plane-space theory into the common geometric foundation: its current
placement near the Part B flatness functional conceals its Chapter 5 use.
